\clearpage\thispagestyle{empty}
\pdfbookmark[0]{Quellen zur Anzeige über die Polhöhe von 1808}{anzeige1808-sources}
\begin{center}\textbf{Quellen zur Anzeige über die Polhöhe von 1808}\end{center}
\noindent Diese KI-Übertragung des deutschen Originals in heutiges Deutsch folgt der Anzeige in Gauß’ \textit{Werken}, Band~VI (Göttingen, 1874), S.~68--69. Die Anzeige erschien in den \textit{Göttingischen gelehrten Anzeigen}, 195.~Stück, 5.~Dezember 1808, S.~1945--1947. \href{https://manifests.sub.uni-goettingen.de/iiif/image/gdz:PPN319721507_1808:00001951/full/1200,/0/default.jpg}{Erstdruck}. \href{https://gdz.sub.uni-goettingen.de/id/PPN236007467}{Historischer deutscher Text in den Werken}.
\par\smallskip\noindent Die Anzeige bespricht Gauß’ \textit{Methodus peculiaris elevationem poli determinandi}, das in diesem Band weiter vorn wiedergegebene Vorlesungsprogramm. Zur heutigen Terminologie von Sternzeit und Polhöhe siehe auch Siegfried Krauter, \textit{Astronomie} (Januar 2012), S.~24--26.
\par\smallskip\noindent Übertragung und Satz wurden mit KI erstellt.
\clearpage
